%!TEX root = ../main.tex

\chapter{Огляд існуючих рішень з автоматичного розпізнавання музичних інструментів}
\label{chap:review}

Автоматичне розпізнавання музичних інструментів є однією з базових задач
музичного інформаційного пошуку. Вона полягає у визначенні за звуковим
сигналом множини інструментів, що звучать у записі. У цьому розділі
розглянуто фізичну природу тембру, на якій ґрунтується розпізнавання,
типовий спосіб часово-частотного представлення звукового сигналу, основні
класи існуючих методів (від класичних методів на основі ручно
сконструйованих ознак до згорткових нейронних мереж) разом із загальними
прийомами їх навчання та застосування, способи врахування гармонічної
будови звуку, методи розширення навчальних даних, комерційні застосунки, а
також набори даних і методики оцінювання. На основі огляду визначено
недоліки існуючих рішень, які обґрунтовують напрям цієї роботи.

% ------------------------------------------------------------------
\section{Постановка задачі розпізнавання музичних інструментів}
\label{sec:task}

Залежно від характеру вхідних даних розрізняють кілька варіантів задачі.
Найпростішим є розпізнавання \textit{ізольованих нот}, записаних окремо для
кожного інструмента в контрольованих умовах. Складнішим є розпізнавання
\textit{сольних фрагментів}, де інструмент виконує музичну фразу з
характерними прийомами гри. Найбільш наближеним до практики є
розпізнавання інструментів у \textit{багатоінструментних записах}, де
кілька джерел звучать одночасно, їхні спектри перекриваються, а до сигналу
додаються ударні, ефекти обробки та шуми запису~\cite{essid2006,fuhrmann2012}.

У багатоінструментному випадку задача природно формулюється як
\textit{багатоміткова класифікація}: кожному запису відповідає довільна
підмножина з $K$ класів інструментів, а не один клас~\cite{tsoumakas2007}.
Часто розглядають її різновид: розпізнавання \textit{переважаючих}
інструментів, тобто тих, що чітко виділяються у звучанні, без вимоги
знайти всі інструменти суміші~\cite{irmas2012}. Особливістю найпоширенішого
набору даних IRMAS є те, що навчальні фрагменти позначено одним
переважаючим інструментом, а тестові одним або кількома. Отже, модель
навчається на одноміткових даних, а оцінюється на багатоміткових, що є
окремим джерелом складності.

% ------------------------------------------------------------------
\section{Тембр і його акустичні ознаки}
\label{sec:timbre}

Тембр є властивістю звуку, яка дає змогу розрізнити два звуки
однакової висоти та гучності. З фізичного погляду тембр визначається
кількома групами характеристик~\cite{muller2015,herrera2003}.

\textit{Гармонічна будова.} Звук більшості мелодичних інструментів є
квазіперіодичним і містить основний тон частоти $f_0$ та обертони на
частотах, близьких до $h f_0$, $h = 2, 3, \ldots$ Співвідношення амплітуд
обертонів характерне для інструмента: так, у звуці кларнета в нижньому
регістрі парні обертони значно слабші за непарні.

\textit{Спектральна обвідна} є плавною функцією, що описує розподіл
енергії за частотою незалежно від конкретної ноти. Вона формується
резонансами корпусу інструмента чи голосового тракту (формантами).

\textit{Часова обвідна} описує зміну амплітуди звуку в часі: швидкість
атаки, наявність сталої фази та характер згасання. Наприклад, фортепіано та
гітара мають різку атаку і поступове згасання, тоді як орган і смичкові
інструменти можуть звучати рівно протягом тривалого часу.

\textit{Шумові складові} (звуки дихання у духових інструментах, тертя
смичка, удар молоточка) також несуть інформацію про джерело звуку.

Ці характеристики є основою як для ручного конструювання ознак у
класичних методах, так і для обґрунтування архітектурних рішень у
нейромережевих моделях.

% ------------------------------------------------------------------
\section{Часово-частотне представлення звукового сигналу}
\label{sec:tfgen}

Більшість методів розпізнавання інструментів, як класичних, так і
нейромережевих, працюють не безпосередньо зі звуковим сигналом, а з його
часово-частотним представленням~\cite{muller2015,choi2017}. Нижче
описано типовий спосіб його обчислення; наведені значення параметрів
поширені в роботах з розпізнавання інструментів~\cite{han2017}.

\textit{Попередня обробка.} Стереосигнал перетворюють на одноканальний
усередненням каналів, передискретизують до частоти $f_s$ (типово
22\,050~Гц) і нормують за піковою амплітудою~\cite{han2017}:
\begin{equation}
    s(n) \leftarrow \frac{s(n)}{\max_n |s(n)|}.
\end{equation}

\textit{Коротковіконне перетворення Фур'є}~\cite{oppenheim2010,muller2015}.
Сигнал ділять на кадри довжиною $N_w$ відліків з кроком $H$ і для кожного
кадру, помноженого на віконну функцію $w(n)$, обчислюють спектр
\begin{equation}
    X(m, k) = \sum_{n=0}^{N_w-1} s(n + mH)\, w(n)\, e^{-2\pi j k n / N_w},
    \quad k = 0, \ldots, N_w/2,
\end{equation}
де $m$ є номером кадру. При $f_s = 22050$~Гц, $N_w = 1024$ і $H = 512$
кадр триває близько 46~мс, а крок становить близько 23~мс.

\textit{Мел-спектрограма.} Мел-шкала пов'язує частоту $f$ у герцах із
суб'єктивно сприйнятою висотою~\cite{stevens1937}; використовується
поширене наближення~\cite{oshaughnessy1987}:
\begin{equation}
    \mu(f) = 2595 \log_{10}\!\left(1 + \frac{f}{700}\right).
    \label{eq:mel}
\end{equation}
Інтервал $[0;\, f_s/2]$ ділять на $F + 1$ рівних відрізків на мел-шкалі,
і центр $b$-ї смуги
\begin{equation}
    \mu_b = b\, \Delta\mu, \quad \Delta\mu = \frac{\mu(f_s/2)}{F+1},
    \quad b = 1, \ldots, F.
    \label{eq:melcentres}
\end{equation}
Банк трикутних фільтрів $\Phi \in \mathbb{R}^{F \times (N_w/2+1)}$,
запропонований для обчислення мел-частотних ознак~\cite{davis1980},
підсумовує енергію спектра в кожній смузі, а логарифмічна мел-спектрограма
\begin{equation}
    M(b, m) = \ln\!\left( \sum_k \Phi_{bk}\, |X(m,k)|^2 + \varepsilon \right),
    \label{eq:logmel}
\end{equation}
де мала стала $\varepsilon$ (наприклад, $10^{-9}$) запобігає логарифму
нуля.

\textit{Часові похідні.} Для відображення часової обвідної, як і в
класичних динамічних ознаках мовних сигналів~\cite{furui1986},
обчислюють першу та другу центральні різниці вздовж осі часу:
\begin{equation}
    \Delta M(b,m) = \frac{M(b,m+1) - M(b,m-1)}{2}, \qquad
    \Delta^2 M = \Delta(\Delta M).
\end{equation}

Логарифмічна мел-спектрограма, іноді разом із похідними, є найпоширенішим
входом згорткових моделей розпізнавання
інструментів~\cite{han2017,pons2017}.

% ------------------------------------------------------------------
\section{Методи на основі ручно сконструйованих ознак}
\label{sec:classical}

Перші системи розпізнавання інструментів ґрунтувалися на обчисленні
набору ознак, що описують тембр, і подальшій класифікації за допомогою
статистичних моделей~\cite{herrera2003}. Найпоширенішими ознаками є
мел-частотні кепстральні коефіцієнти, що компактно описують спектральну
обвідну~\cite{davis1980}, спектральні ознаки (центроїд, розкид, нахил,
спад спектра), ознаки часової обвідної (тривалість атаки, швидкість
згасання) та ознаки гармонічності. Великий перелік таких ознак
систематизовано в роботі~\cite{peeters2004}.

Як класифікатори використовували суміші гауссових розподілів, метод
опорних векторів та метод $k$ найближчих сусідів~\cite{marques1999,
agostini2003}. Для ізольованих нот і сольних фрагментів такі системи
досягали високої точності~\cite{eronen2000}, проте в багатоінструментних
записах їхня якість різко знижується, оскільки ознаки, обчислені за всім
сигналом, змішують внески кількох джерел. Для подолання цієї проблеми
запропоновано ієрархічні класифікатори на основі автоматично побудованих
таксономій ансамблів~\cite{essid2006}, зважування ознак з урахуванням
ступеня їх перекриття~\cite{kitahara2007}, а також попереднє виділення
переважаючого джерела методами розділення звуку~\cite{irmas2012}.

Саме в роботі~\cite{irmas2012} разом із набором даних IRMAS запропоновано
базовий метод на основі ознак та методу опорних векторів, що досягає
макроусередненої F1-міри $0{,}432$ і мікроусередненої $0{,}503$ на
тестовій частині набору~\cite{pons2017}. Загальним недоліком методів цієї
групи є залежність від експертного вибору ознак, які не завжди
відображають відмінності між близькими за звучанням інструментами.

% ------------------------------------------------------------------
\section{Згорткові нейронні мережі}
\label{sec:cnn}

Згорткові нейронні мережі навчаються виділяти ознаки безпосередньо з
часово-частотного представлення сигналу~\cite{lecun2015,goodfellow2016}.
Як вхід зазвичай використовують логарифмічну мел-спектрограму, тобто
двовимірну матрицю «частота $\times$ час», яку мережа обробляє подібно до
зображення.

\subsection{Мережі загального призначення}

Найвідомішим рішенням для набору IRMAS є модель
Хана, Кіма та Лі~\cite{han2017}. Вона складається з чотирьох блоків по дві
згортки $3 \times 3$ з подальшим підвибором, глобального вибору максимуму
та двох повнозв'язних шарів і має близько 1,45~млн параметрів. Модель
навчається на фрагментах тривалістю 1~с, а під час розпізнавання
тестовий запис обробляється ковзним вікном, оцінки вікон підсумовуються та
нормуються на максимум, після чого застосовується поріг. Автори
повідомляють макроусереднену F1-міру $0{,}503$ та мікроусереднену
$0{,}602$. Водночас, згідно з текстом роботи, тривалість вікна, спосіб
агрегації та поріг обирались за найкращим значенням F1-міри, а окремого
опису перевірочної вибірки для цього вибору не наведено, що робить
результат дещо оптимістичним.

Інші роботи досліджували навчання безпосередньо на звуковому
сигналі~\cite{li2015}, спіральне представлення висоти тону, що враховує
октавну подібність~\cite{lostanlen2016}, спільне розпізнавання
інструментів і висоти тону на рівні окремих кадрів~\cite{hung2018}, а також
механізми уваги для розпізнавання за слабкими мітками на наборі
OpenMIC-2018~\cite{gururani2019,openmic2018}.

\subsection{Мережі з музично обґрунтованими фільтрами}

Пояснення роботи мереж загального призначення ускладнене тим, що квадратні
фільтри не мають очевидного акустичного змісту. Понс та
співавтори~\cite{pons2017} запропонували проєктувати форму фільтрів першого
шару відповідно до будови тембру. \textit{Вертикальні} фільтри охоплюють
велику частину частотної осі в межах одного або кількох кадрів і
відображають спектральну обвідну. Фільтри малої висоти та більшої
ширини відображають часову обвідну. У запропонованій ними моделі для
IRMAS перший шар містить фільтри висотою $5$ і $80$ смуг (із $96$) різної
ширини, після нього виконується підвибір майже за всією частотною віссю,
далі два згорткові шари та повнозв'язний шар. Модель має близько
743~тис. параметрів і досягає макроусередненої F1-міри $0{,}516$ та
мікроусередненої $0{,}589$, тобто кращого за модель~\cite{han2017}
значення макроусередненої міри за вдвічі меншої кількості параметрів. Як
і в роботі~\cite{han2017}, поріг рішення ($0{,}2$) та протокол вибору
моделі детально не обґрунтовано, а кількість незалежних запусків не
вказано.

Висоту таких фільтрів задають у частотних смугах мел-спектрограми. Як буде
показано в розділі~\ref{chap:math}, через нерівномірність мел-шкали
фільтр сталої висоти охоплює різні музичні інтервали на різних частотах,
що суперечить ідеї фільтра, який описує тембр незалежно від висоти ноти.

\subsection{Попередньо навчені моделі}

Окремий напрям становлять моделі, попередньо навчені на великих наборах
загальних звуків, зокрема AudioSet~\cite{gemmeke2017,hershey2017,kong2020},
з подальшим донавчанням на цільовій задачі. Такі моделі зазвичай мають
десятки мільйонів параметрів і досягають вищої якості, проте потребують
значних обчислювальних ресурсів, а їхні результати важко зіставляти з
моделями, навченими лише на цільовому наборі. Тому в цій роботі
розглядаються лише моделі, які навчаються з нуля на наборі IRMAS без
залучення зовнішніх даних.

\subsection{Навчання для багатоміткової класифікації}
\label{sec:traingen}

Для багатоміткової постановки вихідний шар мережі містить $K$ незалежних
сигмоїдних виходів $p_k = \sigma(z_k) = 1/(1 + e^{-z_k})$, кожен з яких
оцінює ймовірність присутності відповідного
інструмента~\cite{tsoumakas2007,goodfellow2016}. Параметри мережі
$\theta$ визначають мінімізацією бінарної перехресної ентропії,
усередненої за класами та прикладами пакета з $B$ прикладів:
\begin{equation}
    \mathcal{L}(\theta) = -\frac{1}{BK} \sum_{i=1}^{B} \sum_{k=1}^{K}
    \Big[ y_{ik} \ln p_{ik} + (1 - y_{ik}) \ln(1 - p_{ik}) \Big].
    \label{eq:bce}
\end{equation}
Для мінімізації поширено використовують метод Adam~\cite{kingma2015}, а
крок навчання поступово зменшують, зокрема за косинусним
законом~\cite{loshchilov2017}:
\begin{equation}
    \eta_e = \eta_{\min} + \frac{1}{2}(\eta_0 - \eta_{\min})
    \left(1 + \cos\frac{\pi e}{E}\right),
    \label{eq:cosine}
\end{equation}
де $e$ позначає номер епохи, $E$ задає кількість епох, а $\eta_0$ і
$\eta_{\min}$ є початковим і кінцевим кроком. Для стабілізації навчання
застосовують пакетну нормалізацію~\cite{ioffe2015}, а для регуляризації
проріджування, тобто випадкове вимикання каналів і нейронів під час
навчання~\cite{srivastava2014}.

\subsection{Розпізнавання записів довільної тривалості}
\label{sec:inference}

Моделі навчають на фрагментах фіксованої тривалості (як правило, 1--3~с),
тоді як записи, що аналізуються, довші. У роботі~\cite{han2017}
запропоновано обробляти запис ковзним вікном довжиною $T$ кадрів з кроком
$T/2$, обчислювати для вікон $w = 1, \ldots, W$ ймовірності
$\mathbf{p}^{(w)}$ та агрегувати їх підсумовуванням з нормуванням на
максимум:
\begin{equation}
    S_k = \sum_{w=1}^{W} p^{(w)}_k, \qquad
    \bar{S}_k = \frac{S_k}{\max_{j} S_j}.
    \label{eq:s2}
\end{equation}
Нормування гарантує, що найвпевненіше розпізнаний інструмент отримує
оцінку 1. Остаточне рішення приймається за порогом $\tau$:
\begin{equation}
    \hat{y}_k =
    \begin{cases}
        1, & \bar{S}_k \geq \tau, \\
        0, & \bar{S}_k < \tau,
    \end{cases}
    \qquad k = 1, \ldots, K.
    \label{eq:decision}
\end{equation}
Автори~\cite{han2017} порівнювали також усереднення оцінок без нормування
й обрали нормоване підсумовування з порогом $0{,}5$; у
роботі~\cite{pons2017} оцінки усереднюються, а поріг становить $0{,}2$.

% ------------------------------------------------------------------
\section{Гармонічні та частотно-інваріантні представлення}
\label{sec:harmonic}

Мел-спектрограма має нерівномірну частотну роздільність: смуги рівномірно
розташовані на психоакустичній мел-шкалі~\cite{stevens1937}, яка є майже
лінійною нижче 1~кГц і майже логарифмічною вище. Альтернативою є
перетворення зі сталою добротністю~\cite{brown1991}, у якому кожна октава
містить однакову кількість частотних смуг. На такій осі зміна висоти ноти
відповідає зсуву спектра на сталу кількість смуг, а відстані між
обертонами не залежать від висоти основного тону.

Цю властивість використовують для явного врахування гармонічної будови
звуку. У гармонічному представленні~\cite{bittner2017} кілька копій
спектра, розтягнутих у $h$ разів по частоті, накладаються як канали входу
мережі, так що в одній точці згортка «бачить» енергію основного тону та
обертонів. У роботі~\cite{won2020} запропоновано навчувані гармонічні
фільтри. Інший підхід полягає в розділенні сигналу на гармонічну та ударну
складові методом медіанної фільтрації~\cite{fitzgerald2010}, що дає змогу
зменшити вплив ударних інструментів, які в наборі IRMAS не розмічаються.
Ці представлення здебільшого досліджувались для оцінювання висоти тону та
загальної класифікації звуків, а їхня ефективність для розпізнавання
інструментів у багатоінструментних записах залишається недостатньо
вивченою.

% ------------------------------------------------------------------
\section{Методи розширення навчальних даних}
\label{sec:augmentation}

Набори даних для розпізнавання інструментів невеликі за обсягом, тому
поширеним прийомом є штучне розширення навчальної вибірки. Для звукових
даних застосовують зміну висоти та темпу, додавання шуму, маскування
частотних смуг і часових інтервалів~\cite{park2019,schluter2015}.
Метод змішування прикладів~\cite{zhang2018} формує лінійні комбінації
пар прикладів і їхніх міток, проте для звуку така комбінація логарифмічних
спектрограм не відповідає фізичному додаванню сигналів, а дробові мітки
не відповідають багатомітковій постановці.

Для задачі, в якій навчальні дані містять один переважаючий інструмент, а
тестові кілька, природним є створення штучних
багатоінструментних сумішей з навчальних записів. Такий підхід
досліджувався в роботі~\cite{kratimenos2021}, де показано його користь для
класифікації інструментів у поліфонічній музиці. Водночас питання про те,
як синтетичні суміші взаємодіють з архітектурними рішеннями, зокрема з
гармонічними представленнями, залишається відкритим.

% ------------------------------------------------------------------
\section{Комерційні застосунки}
\label{sec:commercial}

Автоматичне розпізнавання інструментів уже використовується в
комерційних продуктах, як правило, як частина ширших функцій. Сервіси
автоматичного тегування музичних каталогів, зокрема Cyanite~\cite{cyanite},
визначають для записів набір ознак, серед яких жанр, настрій, темп та
інструментальний склад; такі ознаки використовують звукозаписні компанії,
видавці та бібліотеки музики для пошуку й добору творів. Подібні
технології автоматичного тегування та пошуку розвивала компанія Musiio,
яку у 2022~р. придбав сервіс SoundCloud для покращення пошуку
музики~\cite{musiio2022}.

У програмах для зведення звуку розпізнавання інструмента використовується
для автоматичного налаштування обробки: модуль Track Assistant засобу
iZotope Neutron визначає тип інструмента на доріжці (вокал, гітара, бас,
ударні, фортепіано) і пропонує відповідні налаштування~\cite{izotope}.
Близькою задачею є розділення готового запису на складові окремих
інструментів, реалізоване, зокрема, у функції Stem Splitter програми
Logic Pro, яка виділяє вокал, ударні, бас, гітару та
фортепіано~\cite{logicpro2025}.

Спільною рисою комерційних рішень є закритість: методи, навчальні дані
та показники якості, як правило, не публікуються, а набір розпізнаваних
категорій обмежений. Частина рішень працює з окремими доріжками, на яких
звучить один інструмент, а не з готовими багатоінструментними записами.
Тому для порівняння методів використовують відкриті набори даних і
опубліковані результати, розглянуті нижче.

% ------------------------------------------------------------------
\section{Набори даних і методики оцінювання}
\label{sec:datasets}

\textit{IRMAS}~\cite{irmas2012} є стандартним набором даних для
розпізнавання переважаючих інструментів. Він містить 11 класів
(віолончель, кларнет, флейта, акустична гітара, електрогітара, орган,
фортепіано, саксофон, труба, скрипка, голос). Навчальна частина складається
з 6705 фрагментів тривалістю 3~с, вирізаних з понад 2000 записів і
позначених одним переважаючим інструментом; для частини фрагментів у назві
файлу вказано жанр і наявність ударних. Тестова частина містить 2874
фрагменти тривалістю 5--20~с з одним або кількома позначеними
інструментами. Детальніше набір описано в розділі~\ref{sec:data}.

\textit{OpenMIC-2018}~\cite{openmic2018} містить близько 20~тис.
десятисекундних фрагментів з 20 класами та частковою (слабкою)
розміткою. \textit{NSynth}~\cite{engel2017} містить понад 300~тис.
ізольованих нот і придатний переважно для задачі розпізнавання окремих
звуків.

Для оцінювання багатоміткового розпізнавання використовують точність,
повноту та F1-міру для кожного класу, а також їх макроусереднення (середнє
за класами) та мікроусереднення (обчислення за об'єднанням усіх
рішень)~\cite{sokolova2009}. Макроусереднена міра однаково враховує всі
класи, тому чутлива до якості розпізнавання рідкісних інструментів.

Аналіз опублікованих робіт показує, що методики оцінювання на IRMAS
суттєво різняться і часто описані неповно: не завжди вказано спосіб
формування перевірочної вибірки, правило вибору порогу та кількість
незалежних запусків. Окремою проблемою є те, що навчальні фрагменти IRMAS
вирізано з довших записів по кілька з кожного, тому випадковий поділ
фрагментів на навчальну та перевірочну вибірки призводить до потрапляння
частин одного запису в обидві вибірки і завищеної оцінки якості на
перевірочних даних.

% ------------------------------------------------------------------
\section{Висновки до розділу}
\label{sec:ch1_conclusions}

У розділі розглянуто існуючі рішення задачі автоматичного розпізнавання
музичних інструментів. Узагальнене порівняння основних груп методів
наведено в таблиці~\ref{tab:review}.

\begin{table}[H]
\centering
\caption{Порівняння груп методів розпізнавання музичних інструментів}
\label{tab:review}
\small
\begin{tabular}{|p{3.4cm}|p{4.6cm}|p{4.6cm}|p{2.6cm}|}
\hline
\textbf{Група методів} & \textbf{Переваги} & \textbf{Недоліки} &
\textbf{Макро-F1 на IRMAS} \\
\hline
Ручно сконструйовані ознаки та класифікатори~\cite{irmas2012} &
Інтерпретованість, невеликі вимоги до обчислень &
Залежність від вибору ознак, низька якість у сумішах &
0,432 \\
\hline
Згорткові мережі загального призначення~\cite{han2017} &
Автоматичне виділення ознак &
Велика кількість параметрів, відсутність акустичного змісту фільтрів &
0,503 \\
\hline
Мережі з музично обґрунтованими фільтрами~\cite{pons2017} &
Менша кількість параметрів, інтерпретованість фільтрів &
Висота фільтрів задана в мел-смугах, а не в музичних інтервалах &
0,516 \\
\hline
Попередньо навчені моделі~\cite{kong2020} &
Висока якість завдяки великим зовнішнім даним &
Десятки мільйонів параметрів, складність порівняння &
- \\
\hline
Комерційні сервіси~\cite{cyanite,izotope} &
Готові до використання, працюють з великими каталогами &
Закриті методи та показники якості, обмежений набір категорій &
- \\
\hline
\end{tabular}
\end{table}

За результатами огляду виявлено такі недоліки існуючих рішень:
\begin{enumerate}
    \item вертикальні фільтри, що мають описувати тембр, задаються сталою
    висотою в смугах мел-спектрограми і через нерівномірність мел-шкали
    охоплюють різні музичні інтервали на різних частотах;
    \item моделі навчаються на записах з одним переважаючим інструментом, а
    застосовуються до багатоінструментних, причому цю розбіжність, як
    правило, не компенсують на етапі навчання;
    \item вплив явних гармонічних представлень на розпізнавання
    інструментів у сумішах, особливо в поєднанні з розширенням даних,
    недостатньо досліджено;
    \item методики оцінювання на наборі IRMAS неоднорідні: поріг рішення
    іноді обирають за результатами на тестових даних, кількість запусків
    не вказують, а випадковий поділ фрагментів спричиняє перетин
    навчальної та перевірочної вибірок.
\end{enumerate}

Отже, у роботі доцільно розробити модель згорткової нейронної мережі, у
якій тембральні фільтри охоплюють сталі музичні інтервали, навчати її на
синтетичних багатоінструментних сумішах, дослідити роль гармонічних
представлень та оцінювати результати за відтворюваною методикою без
використання тестових даних для вибору налаштувань.
